% Abstract (hand-written)
\chapter*{Abstract}
\addcontentsline{toc}{chapter}{Abstract}
Asking an AI about something on the screen normally takes seven manual steps: take a
screenshot, crop it, open a chatbot, paste the image, retype the context, wait, and switch
back. IntelliPointer AI reduces this to one shortcut. A Windows application freezes the
screen, lets the user mark one or more regions, reads the text on the PC, and shows a streamed
answer from a vision-capable model in a small floating window, with follow-up questions that
reuse the same screenshot.

A model needs a secret key, and a key inside a desktop program can be extracted. The project
therefore places a stateless HTTPS API in the cloud (Render or Google Cloud Run). Firebase
Authentication identifies each user with short-lived signed tokens, and Cloud Firestore holds
a per-user daily quota that is updated in an atomic transaction, so concurrent requests cannot
exceed it. The model key never leaves the server, nothing the user captures or asks is stored
or logged, and every control is covered by an automated test, including integration tests
against the real Firebase emulators. On the test machine the overlay appears in under 40 ms
and a finished crop is ready in about 0.2 to 0.5 s.

\bigskip
\noindent\textbf{Keywords:} cloud computing, Firebase, Firestore, serverless, Server-Sent
Events, authentication, quota, desktop application, vision model.
\clearpage
